\documentclass[aspectratio=169]{beamer}

% =====================================================================
% ART DECK 02 -- no limits. The most beautiful version I can build.
% True soft glow via radial fadings, screen blend modes so light adds,
% layered nebulae, bokeh fields, lens-flare stars, volumetric rays,
% neon triple-stroke wires, aurora curtains. Full-bleed 16x9 canvas.
% Two pdflatex passes (remember picture/overlay).
% =====================================================================

\usepackage{tikz}
\usepackage[T1]{fontenc}
\usetikzlibrary{shadings, fadings, shapes.geometric, shapes.misc,
  decorations.pathmorphing, decorations.pathreplacing, calc, backgrounds,
  positioning, arrows.meta, shapes.symbols}

\usetheme{default}
\setbeamertemplate{navigation symbols}{}
\setbeamercolor{background canvas}{bg=black}
\setbeamertemplate{footline}{}

% ---- rich neon palette ----------------------------------------------------
\definecolor{indigo}{HTML}{241A52}
\definecolor{deepblue}{HTML}{0C1733}
\definecolor{midnight}{HTML}{070C1C}
\definecolor{ncyan}{HTML}{38E8FF}
\definecolor{nteal}{HTML}{19E6CE}
\definecolor{nmag}{HTML}{FF49B0}
\definecolor{npink}{HTML}{FF77D6}
\definecolor{nlime}{HTML}{B9FF45}
\definecolor{namber}{HTML}{FFC24A}
\definecolor{gold}{HTML}{FFD56B}
\definecolor{eblue}{HTML}{4F7BFF}
\definecolor{royal}{HTML}{6A5BFF}
\definecolor{nviolet}{HTML}{A05CFF}
\definecolor{ngreen}{HTML}{3CE89A}
\definecolor{nred}{HTML}{FF5468}
\definecolor{faint}{HTML}{6E7E92}
\definecolor{paper}{HTML}{EAF2FB}

% ---- radial / axial fadings for true glow ---------------------------------
\tikzfading[name=rglow,  inner color=transparent!0,  outer color=transparent!100]
\tikzfading[name=rcore,  inner color=transparent!0,  outer color=transparent!100]
\tikzfading[name=fadeup, top color=transparent!0,    bottom color=transparent!100]
\tikzfading[name=fadedn, top color=transparent!100,  bottom color=transparent!0]

\pgfmathsetseed{777}

% ---- full-bleed art frame: (0,0)=SW .. (16,9)=NE --------------------------
\newcommand{\art}[1]{%
  \begin{frame}[plain]
  \begin{tikzpicture}[remember picture, overlay,
      shift={(current page.south west)}, x=1cm, y=1cm]
    #1
  \end{tikzpicture}
  \end{frame}}

% soft round halo: {(c)}{r}{color}{opacity}
\newcommand{\halo}[4]{\fill[#3, opacity=#4, path fading=rglow] #1 circle (#2);}

% a luminous star with lens-flare cross: {(c)}{r}{color}
\newcommand{\spark}[3]{%
  \fill[#3, opacity=0.55, path fading=rglow] #1 circle ({#2*7});
  \fill[#3, opacity=0.85, path fading=rglow] #1 circle ({#2*2.8});
  \begin{scope}[blend mode=screen]
    \fill[white, opacity=0.55, path fading=rglow] #1 ellipse ({#2*9} and {#2*0.45});
    \fill[white, opacity=0.55, path fading=rglow] #1 ellipse ({#2*0.45} and {#2*9});
  \end{scope}
  \fill[white] #1 circle (#2);}

% neon triple-stroke straight wire: {(a)}{(b)}{color}
\newcommand{\neon}[3]{%
  \draw[#3, line width=7pt, opacity=0.10] #1 -- #2;
  \draw[#3, line width=3pt, opacity=0.30] #1 -- #2;
  \draw[white, line width=0.7pt, opacity=0.80] #1 -- #2;}

% dreamy bokeh field under screen blend: {count}
\newcommand{\bokeh}[1]{%
  \begin{scope}[blend mode=screen]
  \foreach \i in {1,...,#1}{%
     \pgfmathsetmacro{\bx}{rnd*16}\pgfmathsetmacro{\by}{rnd*9}%
     \pgfmathsetmacro{\br}{0.05+rnd*0.50}\pgfmathsetmacro{\bo}{0.10+rnd*0.38}%
     \pgfmathtruncatemacro{\ci}{rnd*4.999}%
     \ifcase\ci \def\bc{ncyan}\or\def\bc{nviolet}\or\def\bc{nmag}\or
        \def\bc{eblue}\or\def\bc{nteal}\fi
     \fill[\bc, opacity=\bo, path fading=rglow] (\bx,\by) circle (\br);}
  \end{scope}}

% sharp star dust: {count}
\newcommand{\dust}[1]{%
  \foreach \i in {1,...,#1}{%
     \pgfmathsetmacro{\sx}{rnd*16}\pgfmathsetmacro{\sy}{rnd*9}%
     \pgfmathsetmacro{\sr}{0.005+rnd*0.020}\pgfmathsetmacro{\so}{0.10+rnd*0.55}%
     \fill[white, opacity=\so] (\sx,\sy) circle (\sr);}}

% volumetric light rays from a point: {(c)}{color}{count}
\newcommand{\rays}[3]{%
  \begin{scope}[blend mode=screen]
  \foreach \k in {1,...,#3}{%
     \pgfmathsetmacro{\ang}{rnd*360}\pgfmathsetmacro{\len}{4.5+rnd*7}%
     \pgfmathsetmacro{\wd}{0.4+rnd*1.4}%
     \fill[#2, opacity=0.045, path fading=rglow]
        #1 -- ($#1+(\ang-\wd:\len)$) -- ($#1+(\ang+\wd:\len)$) -- cycle;}
  \end{scope}}

% glowing title: {(c)}{halo color}{node text}
\newcommand{\glowtext}[3]{%
  \begin{scope}[blend mode=screen]
     \fill[#2, opacity=0.55, path fading=rglow] #1 ellipse (5.2 and 1.5);
  \end{scope}
  \node at #1 {#3};}

\begin{document}

% =====================================================================
% 1. COVER -- a deep-space nebula with a luminous network
% =====================================================================
\art{
  \fill[midnight] (0,0) rectangle (16,9);
  % nebula clouds
  \begin{scope}[blend mode=screen]
     \fill[royal,  opacity=0.40, path fading=rglow] (5,6) circle (5.5);
     \fill[nmag,   opacity=0.30, path fading=rglow] (11,6.5) circle (4.5);
     \fill[ncyan,  opacity=0.26, path fading=rglow] (9,3) circle (5);
     \fill[nviolet,opacity=0.34, path fading=rglow] (3,3) circle (4.5);
     \fill[indigo, opacity=0.55, path fading=rglow] (8,4.5) circle (8);
  \end{scope}
  \dust{220}\bokeh{60}
  % constellation network
  \coordinate (n1) at (3.6,7.0); \coordinate (n2) at (6.6,7.9);
  \coordinate (n3) at (10.2,7.1);\coordinate (n4) at (2.8,4.8);
  \coordinate (n5) at (6.2,5.3); \coordinate (n6) at (10.9,5.1);
  \coordinate (n7) at (13.0,6.3);\coordinate (n8) at (8.3,3.6);
  \coordinate (n9) at (4.9,3.1);
  \foreach \a/\b/\c in {n1/n2/ncyan,n2/n3/nviolet,n1/n4/ncyan,n4/n5/eblue,
     n5/n6/nmag,n3/n6/nviolet,n2/n5/ncyan,n6/n7/nmag,n5/n8/eblue,
     n8/n9/ncyan,n4/n9/eblue,n9/n5/nviolet}{\neon{(\a)}{(\b)}{\c}}
  \foreach \n/\c in {n1/ncyan,n2/nviolet,n3/nmag,n4/eblue,n5/nteal,
     n6/npink,n7/gold,n8/eblue,n9/nviolet}{\spark{(\n)}{0.07}{\c}}
  % title
  \begin{scope}[blend mode=screen]
     \fill[ncyan, opacity=0.40, path fading=rglow] (5.6,1.75) ellipse (5.6 and 1.5);
  \end{scope}
  \node[anchor=west, text=white] at (1.3,2.05)
    {\fontsize{46}{50}\selectfont\textbf{Rotor \textcolor{ncyan}{in} Rust}};
  \node[anchor=west, text=paper, opacity=0.9] at (1.4,1.05)
    {\large Models that ask one question of every input.};
  \node[anchor=west, text=paper, opacity=0.55] at (1.4,0.42)
    {\footnotesize Jasmin Begi\'{c}\; \textcolor{nmag}{$\bullet$}\; Daniel Wassie};
}

% =====================================================================
% 2. THE PROBLEM -- a supernova among countless inputs
% =====================================================================
\art{
  \fill[midnight] (0,0) rectangle (16,9);
  \begin{scope}[blend mode=screen]
     \fill[indigo, opacity=0.5, path fading=rglow] (8,4.5) circle (9);
     \fill[nmag,   opacity=0.5, path fading=rglow] (11.3,5.7) circle (4.5);
  \end{scope}
  % the field of dim inputs
  \foreach \x in {0,...,37}{\foreach \y in {0,...,16}{
     \pgfmathsetmacro{\jx}{\x*0.40+0.6+rnd*0.05}
     \pgfmathsetmacro{\jy}{\y*0.40+1.5+rnd*0.05}
     \fill[eblue!60!midnight, opacity={0.30+rnd*0.30}] (\jx,\jy) circle (0.045);}}
  \dust{120}
  % the one crashing input -- supernova
  \rays{(11.3,5.7)}{npink}{26}
  \spark{(11.3,5.7)}{0.12}{nmag}
  \node[text=white] at (8,8.4) {\Large the one input that crashes \textcolor{npink}{hides} among};
  \glowtext{(8,0.85)}{ncyan}{\fontsize{36}{38}\selectfont\textbf{\textcolor{ncyan}{$2^{64}$} \normalsize\textcolor{paper}{possibilities}}}
}

% =====================================================================
% 3. VERIFICATION -- one luminous beam over a sea of inputs
% =====================================================================
\art{
  \fill[deepblue!85!black] (0,0) rectangle (16,9);
  \begin{scope}[blend mode=screen]
     \fill[royal, opacity=0.4, path fading=rglow] (8,7) circle (7);
  \end{scope}
  % sea of inputs, all lit
  \foreach \x in {0,...,37}{\foreach \y in {0,...,12}{
     \pgfmathsetmacro{\jx}{\x*0.40+0.6}
     \pgfmathsetmacro{\jy}{\y*0.40+1.6}
     \fill[ncyan, opacity={0.20+rnd*0.5}] (\jx,\jy) circle (0.05);}}
  % volumetric beam
  \begin{scope}[blend mode=screen]
     \fill[ncyan, opacity=0.20, path fading=rglow]
        (8,8.4) -- (0.8,1.3) -- (15.2,1.3) -- cycle;
     \fill[white, opacity=0.30, path fading=rglow]
        (8,8.4) -- (5.4,1.3) -- (10.6,1.3) -- cycle;
  \end{scope}
  \spark{(8,8.3)}{0.12}{ncyan}
  \node[text=paper] at (8,0.7)
    {\Large one question \;$\Rightarrow$\; \textcolor{ncyan}{every} input, at once};
}

% =====================================================================
% 4. THE PIPELINE -- a river of light
% =====================================================================
\art{
  \fill[midnight] (0,0) rectangle (16,9);
  \begin{scope}[blend mode=screen]
     \fill[indigo, opacity=0.55, path fading=rglow] (8,4.6) circle (8.5);
     \fill[nteal,  opacity=0.18, path fading=rglow] (13,4.6) circle (4);
     \fill[eblue,  opacity=0.18, path fading=rglow] (3,4.6) circle (4);
  \end{scope}
  \dust{70}
  \def\sy{4.6}
  % glowing flow wire under the stations
  \draw[ncyan, line width=10pt, opacity=0.10] (2.2,\sy) -- (13.4,\sy);
  \draw[nviolet,line width=4pt,  opacity=0.25] (2.2,\sy) -- (13.4,\sy);
  \foreach \x/\lbl/\c in {2.2/binary/eblue, 6.0/model/ncyan,
        9.8/solver/nviolet, 13.4/answer/nlime}{
     \fill[\c, opacity=0.55, path fading=rglow] (\x,\sy) circle (1.7);
     \shade[inner color=white, outer color=\c!55!midnight] (\x,\sy) circle (0.72);
     \draw[\c, line width=1.4pt] (\x,\sy) circle (0.72);
     \node[text=midnight, font=\footnotesize\bfseries] at (\x,\sy) {\lbl};}
  \foreach \a/\b in {2.2/6.0, 6.0/9.8, 9.8/13.4}{
     \begin{scope}[blend mode=screen]
        \fill[white, opacity=0.8, path fading=rglow]
           ({(\a+\b)/2},\sy) circle (0.16);\end{scope}}
  \node[text=paper, opacity=0.92] at (8,8.0)
    {\Large from a \textcolor{eblue}{compiled binary} to a \textcolor{nlime}{readable answer}};
  \node[text=ncyan, font=\footnotesize\bfseries] at (6.0,2.5) {rotor};
  \node[text=nlime, font=\footnotesize\bfseries] at (13.4,2.5){visualizer};
}

% =====================================================================
% 5. SPEED -- a hyperspace comet
% =====================================================================
\art{
  \fill[midnight] (0,0) rectangle (16,9);
  \begin{scope}[blend mode=screen]
     \fill[indigo, opacity=0.6, path fading=rglow] (8,4.5) circle (10);
     \fill[namber, opacity=0.25, path fading=rglow] (13.4,4.9) circle (4.5);
  \end{scope}
  \dust{120}
  % motion streaks
  \begin{scope}[blend mode=screen]
  \foreach \i in {1,...,40}{
     \pgfmathsetmacro{\my}{1.2+rnd*7.2}\pgfmathsetmacro{\mx}{1+rnd*9}
     \pgfmathsetmacro{\ml}{0.4+rnd*2.2}\pgfmathsetmacro{\mo}{0.05+rnd*0.18}
     \fill[ncyan, opacity=\mo, path fading=rglow]
        (\mx,\my) ellipse ({\ml} and 0.02);}
  \end{scope}
  % comet tail
  \begin{scope}[blend mode=screen]
     \fill[namber, opacity=0.55, path fading=rglow]
        (1.5,2.7) .. controls (6,3.4) and (10,4.2) .. (13.4,4.9)
        .. controls (10,3.7) and (6,3.0) .. (1.5,2.5) -- cycle;
     \fill[gold, opacity=0.5, path fading=rglow]
        (4.5,3.3) .. controls (8,3.9) and (11,4.4) .. (13.4,4.9)
        .. controls (11,4.1) and (8,3.6) .. (4.5,3.1) -- cycle;
  \end{scope}
  \rays{(13.5,4.9)}{gold}{22}
  \spark{(13.5,4.9)}{0.16}{namber}
  \node[text=white] at (5.2,7.2) {\fontsize{70}{72}\selectfont\textbf{1000\textcolor{namber}{$\times$}}};
  \node[text=paper, opacity=0.9] at (5.4,5.7) {\large faster model generation};
  \node[text=paper, opacity=0.78] at (4.2,1.3)
    {\footnotesize \textcolor{nred}{139\,s} $\rightarrow$ \textcolor{gold}{0.1\,s}
     \qquad \textcolor{nred}{428\,MB} $\rightarrow$ \textcolor{gold}{20\,MB}};
}

% =====================================================================
% 6. THE COUNTING -- a tower of light vs a single point
% =====================================================================
\art{
  \fill[midnight] (0,0) rectangle (16,9);
  \begin{scope}[blend mode=screen]
     \fill[nred,  opacity=0.45, path fading=rglow] (4.4,4.6) circle (5);
     \fill[ncyan, opacity=0.40, path fading=rglow] (11.6,2.2) circle (3.5);
  \end{scope}
  \dust{90}
  % the towering bar
  \begin{scope}[blend mode=screen]
     \fill[nred, opacity=0.5, path fading=rglow] (4.4,4.6) ellipse (1.7 and 4.2);
  \end{scope}
  \shade[bottom color=nred!70!black, top color=npink] (3.5,0.9) rectangle (5.3,8.3);
  \draw[npink, line width=1pt, opacity=0.6] (3.5,0.9) rectangle (5.3,8.3);
  \node[text=white, rotate=90, font=\footnotesize] at (4.4,4.6)
     {11{,}695{,}232{,}963 comparisons};
  % the single probe
  \spark{(11.6,2.2)}{0.13}{ncyan}
  \node[text=ncyan, font=\footnotesize] at (11.6,3.2) {159{,}018 hash probes};
  \node[text=paper] at (10.9,7.9) {\Large \textcolor{nred}{a list walk}};
  \node[text=paper] at (10.9,6.9) {\Large vs.\ \textcolor{ncyan}{an index}};
}

% =====================================================================
% 7. DEDUP OFF -- a shattering crystal vs a serene gem
% =====================================================================
\art{
  \fill[deepblue!88!black] (0,0) rectangle (16,9);
  \begin{scope}[blend mode=screen]
     \fill[nred,  opacity=0.40, path fading=rglow] (4.3,4.6) circle (4);
     \fill[ncyan, opacity=0.40, path fading=rglow] (11.7,4.6) circle (4);
  \end{scope}
  \dust{70}
  % shattering red crystal + discharge
  \rays{(4.3,4.6)}{nred}{16}
  \shade[inner color=npink, outer color=nred!60!black]
     (4.3,6.5) -- (3.0,4.4) -- (3.8,2.4) -- (4.9,2.7) -- (5.6,4.7) -- cycle;
  \begin{scope}[blend mode=screen]
     \foreach \a/\b in {6.5/4.2, 4.4/3.8, 4.7/4.7, 2.4/4.4}{}
  \end{scope}
  \draw[white, opacity=0.9, line width=1pt] (4.3,6.5) -- (4.0,4.2);
  \draw[white, opacity=0.9, line width=1pt] (3.0,4.4) -- (4.4,3.7);
  \draw[white, opacity=0.9, line width=1pt] (5.6,4.7) -- (4.0,4.2);
  \draw[white, opacity=0.9, line width=1pt] (3.8,2.4) -- (4.6,4.4);
  % a cracked-off shard flying
  \fill[nred!70!black, opacity=0.9] (6.1,5.6) -- (6.6,5.0) -- (6.3,4.5) -- cycle;
  \node[text=white, font=\bfseries] at (4.3,1.4) {C: crashes};
  % serene faceted gem with caustics
  \shade[inner color=white, outer color=eblue!80!black]
     (11.7,6.4) -- (10.3,4.6) -- (11.7,2.7) -- (13.1,4.6) -- cycle;
  \draw[ncyan, opacity=0.7, line width=0.8pt]
     (11.7,6.4) -- (11.7,2.7) (10.3,4.6) -- (13.1,4.6)
     (11.7,6.4) -- (10.3,4.6) -- (11.7,2.7) -- (13.1,4.6) -- cycle;
  \begin{scope}[blend mode=screen]
     \fill[white, opacity=0.7, path fading=rglow] (11.7,5.0) circle (0.5);
  \end{scope}
  \node[text=white, font=\bfseries] at (11.7,1.4) {Rust: unchanged};
  \node[text=paper] at (8,8.3) {\Large turn off the duplicate check\dots};
}

% =====================================================================
% 8. THE BACK-PORT -- twin suns bridged by plasma
% =====================================================================
\art{
  \fill[midnight] (0,0) rectangle (16,9);
  \begin{scope}[blend mode=screen]
     \fill[indigo, opacity=0.55, path fading=rglow] (8,4.7) circle (8);
     \fill[namber, opacity=0.30, path fading=rglow] (3.9,4.7) circle (3.2);
     \fill[ncyan,  opacity=0.30, path fading=rglow] (12.1,4.7) circle (3.2);
  \end{scope}
  \dust{70}
  % twin suns
  \rays{(3.9,4.7)}{gold}{16}\rays{(12.1,4.7)}{ncyan}{16}
  \shade[inner color=white, outer color=namber!70!black] (3.9,4.7) circle (1.05);
  \node[text=midnight, font=\bfseries] at (3.9,4.7) {Rust};
  \shade[inner color=white, outer color=ncyan!70!black] (12.1,4.7) circle (1.05);
  \node[text=midnight, font=\bfseries] at (12.1,4.7) {C};
  % plasma arc (jagged) + glow
  \draw[white, line width=6pt, opacity=0.10] (5.0,4.7) -- (11.0,4.7);
  \draw[nlime, line width=2.5pt, opacity=0.30,
        decorate, decoration={random steps, segment length=5pt, amplitude=2pt}]
        (5.0,4.7) -- (11.0,4.7);
  \draw[white, line width=1pt, opacity=0.85,
        decorate, decoration={random steps, segment length=5pt, amplitude=2pt}]
        (5.0,4.7) -- (11.0,4.7);
  \spark{(8,4.7)}{0.10}{nlime}
  \node[text=paper] at (8,8.0) {\Large the same idea, carried back into C};
  \glowtext{(8,2.2)}{nlime}{\fontsize{40}{42}\selectfont\textbf{\textcolor{white}{95}\textcolor{ncyan}{$\times$}}}
  \node[text=paper, opacity=0.78] at (8,1.2)
    {\footnotesize byte-identical \;\textcolor{nlime}{$\bullet$}\; 36/36 equivalent};
}

% =====================================================================
% 9. SYMBOLIC ARGV -- locks dissolving into light
% =====================================================================
\art{
  \fill[deepblue!85!black] (0,0) rectangle (16,9);
  \begin{scope}[blend mode=screen]
     \fill[royal, opacity=0.5, path fading=rglow] (8,4.6) circle (8);
  \end{scope}
  \dust{80}
  \foreach \i/\open/\c in {0/0/faint,1/0/faint,2/1/ncyan,3/1/nmag,4/1/nlime,5/1/namber}{
     \pgfmathsetmacro{\lx}{2.0+\i*2.35}
     \ifnum\open=1
        \fill[\c, opacity=0.5, path fading=rglow] (\lx,4.7) circle (1.5);
        % bytes streaming upward
        \begin{scope}[blend mode=screen]
        \foreach \k in {1,...,7}{
           \pgfmathsetmacro{\dy}{rnd*2.4}\pgfmathsetmacro{\dx}{(rnd-0.5)*0.7}
           \fill[\c, opacity={0.5-0.15*rnd}, path fading=rglow]
              (\lx+\dx,5.4+\dy) circle (0.07);}
        \end{scope}
     \fi
     \shade[top color=\c!75!deepblue, bottom color=\c!30!deepblue]
        (\lx-0.5,3.5) rectangle (\lx+0.5,4.7);
     \draw[\c, line width=1.5pt] (\lx-0.5,3.5) rectangle (\lx+0.5,4.7);
     \ifnum\open=1
        \draw[\c, line width=1.6pt] (\lx-0.30,4.7) arc (200:20:0.36);
     \else
        \draw[\c, line width=1.6pt] (\lx-0.30,4.7) arc (180:0:0.30);
     \fi
     \fill[white] (\lx,4.1) circle (0.08);}
  \node[text=paper] at (8,8.2) {\Large the command line: \textcolor{faint}{frozen} $\rightarrow$ \textcolor{ncyan}{open}};
  \node[text=paper, opacity=0.85] at (8,2.0) {\large the solver fills the open bytes};
}

% =====================================================================
% 10. THE X / Y SECRET -- twin dials with a radiant detonation
% =====================================================================
\art{
  \fill[midnight] (0,0) rectangle (16,9);
  \begin{scope}[blend mode=screen]
     \fill[nmag, opacity=0.5, path fading=rglow] (8,5.0) circle (5.5);
     \fill[royal,opacity=0.4, path fading=rglow] (8,5.0) circle (8);
  \end{scope}
  \dust{80}
  \foreach \dx/\ch/\c in {5.4/X/ncyan, 10.6/Y/nlime}{
     \fill[\c, opacity=0.45, path fading=rglow] (\dx,5.1) circle (2.1);
     \shade[inner color=\c!35!midnight, outer color=midnight] (\dx,5.1) circle (1.35);
     \draw[\c, line width=2pt] (\dx,5.1) circle (1.35);
     \draw[\c, line width=0.8pt, opacity=0.6] (\dx,5.1) circle (1.05);
     \foreach \t in {0,15,...,345}{\draw[\c,opacity=0.55,line width=1pt]
        ({\dx+1.16*cos(\t)},{5.1+1.16*sin(\t)}) --
        ({\dx+1.35*cos(\t)},{5.1+1.35*sin(\t)});}
     \node[text=white] at (\dx,5.1) {\fontsize{40}{42}\selectfont\textbf{\ch}};}
  % detonation between
  \rays{(8,5.1)}{gold}{30}
  \begin{scope}[blend mode=screen]
  \foreach \t in {0,18,...,342}{
     \pgfmathsetmacro{\rr}{0.7+0.6*rnd}
     \draw[namber,opacity=0.8,line width=1.2pt]
        ({8+0.3*cos(\t)},{5.1+0.3*sin(\t)}) -- ({8+\rr*cos(\t)},{5.1+\rr*sin(\t)});}
  \end{scope}
  \spark{(8,5.1)}{0.10}{namber}
  \node[text=paper] at (8,8.5) {\Large two letters, two arguments --- found by logic};
  \node[text=npink] at (8,1.2) {\large the planted secret, recovered in seconds};
}

% =====================================================================
% 11. 36 / 36 -- a radiant bloom of checks
% =====================================================================
\art{
  \fill[deepblue!85!black] (0,0) rectangle (16,9);
  \begin{scope}[blend mode=screen]
     \fill[ngreen, opacity=0.45, path fading=rglow] (8,5.4) circle (7);
     \fill[nteal,  opacity=0.30, path fading=rglow] (8,5.4) circle (4);
  \end{scope}
  \dust{70}\rays{(8,5.4)}{ngreen}{18}
  \foreach \r in {0,...,3}{\foreach \cc in {0,...,8}{
     \pgfmathsetmacro{\cx}{2.6+\cc*1.35}
     \pgfmathsetmacro{\cy}{7.0-\r*1.05}
     \fill[ngreen, opacity=0.5, path fading=rglow] (\cx,\cy) circle (0.42);
     \draw[white, line width=1.8pt, line cap=round]
        (\cx-0.18,\cy) -- (\cx-0.03,\cy-0.16) -- (\cx+0.24,\cy+0.22);}}
  \glowtext{(8,1.7)}{ngreen}{\fontsize{46}{48}\selectfont\textbf{\textcolor{white}{36} \textcolor{ngreen}{/} \textcolor{white}{36}}}
  \node[text=paper, opacity=0.85] at (8,0.75)
    {\footnotesize same bug, same step --- both tools, every benchmark};
}

% =====================================================================
% 12. THE VISUALIZER -- the model as a luminous neural constellation
% =====================================================================
\art{
  \fill[midnight] (0,0) rectangle (16,9);
  \begin{scope}[blend mode=screen]
     \fill[nred,    opacity=0.40, path fading=rglow] (8,7.3) circle (3.5);
     \fill[nviolet, opacity=0.35, path fading=rglow] (8,4) circle (7);
  \end{scope}
  \dust{90}
  \coordinate (top) at (8,6.7);
  \coordinate (na) at (5.6,5.4); \coordinate (nb) at (10.4,5.4);
  \coordinate (nc) at (3.4,3.6); \coordinate (nd) at (6.8,3.6);
  \coordinate (ne) at (9.2,3.6); \coordinate (nf) at (12.2,3.6);
  \coordinate (ng) at (4.8,1.9); \coordinate (nh) at (7.6,1.9);
  \foreach \a/\b/\c in {top/na/nred, top/nb/nred, na/nc/ncyan, na/nd/nviolet,
        nb/ne/ncyan, nb/nf/nlime, nd/ng/namber, nd/nh/nviolet}{\neon{(\a)}{(\b)}{\c}}
  % bad-state star with rays
  \rays{(8,7.4)}{nred}{20}
  \node[star, star points=8, star point ratio=0.55, draw=nred, fill=nred!75!black,
        line width=1.3pt, minimum size=1.7cm] at (8,7.4) {};
  \node[text=white, font=\tiny\bfseries] at (8,7.4) {VIOLATED};
  \foreach \n/\c in {na/ncyan, nb/nteal, nc/eblue, nd/nviolet, ne/eblue,
        nf/nlime, ng/namber, nh/npink}{\spark{(\n)}{0.085}{\c}}
  \node[text=paper] at (8,0.8) {\Large the unreadable answer, made \textcolor{nlime}{watchable}};
}

% =====================================================================
% 13. THE LESSON -- a luminous crest of trust
% =====================================================================
\art{
  \fill[midnight] (0,0) rectangle (16,9);
  \begin{scope}[blend mode=screen]
     \fill[royal, opacity=0.55, path fading=rglow] (8,5) circle (8);
     \fill[ncyan, opacity=0.30, path fading=rglow] (8,5.2) circle (3.2);
  \end{scope}
  \dust{110}\rays{(8,5.2)}{ncyan}{16}
  % shield
  \shade[inner color=white, outer color=eblue!55!midnight]
     (8,7.8) -- (5.8,6.6) -- (5.8,3.9)
     .. controls (5.8,2.6) and (7.0,2.1) .. (8,1.6)
     .. controls (9.0,2.1) and (10.2,2.6) .. (10.2,3.9)
     -- (10.2,6.6) -- cycle;
  \draw[ncyan, line width=2pt]
     (8,7.8) -- (5.8,6.6) -- (5.8,3.9)
     .. controls (5.8,2.6) and (7.0,2.1) .. (8,1.6)
     .. controls (9.0,2.1) and (10.2,2.6) .. (10.2,3.9)
     -- (10.2,6.6) -- cycle;
  \begin{scope}[blend mode=screen]
     \fill[white, opacity=0.8, path fading=rglow] (7.7,5.2) circle (1.3);
  \end{scope}
  \draw[midnight, line width=3pt, line cap=round, line join=round]
     (7.2,4.8) -- (7.75,3.9) -- (8.9,5.9);
  \node[text=paper] at (8,0.85)
    {\large the hard part is \textcolor{ncyan}{proving it still tells the truth}};
}

% =====================================================================
% 14. FINALE -- aurora over a nebula, the constellation, the close
% =====================================================================
\art{
  \fill[midnight] (0,0) rectangle (16,9);
  % nebula
  \begin{scope}[blend mode=screen]
     \fill[nviolet,opacity=0.40, path fading=rglow] (6,6) circle (6);
     \fill[nmag,   opacity=0.30, path fading=rglow] (11,5) circle (5);
     \fill[ncyan,  opacity=0.26, path fading=rglow] (9,3) circle (5.5);
     \fill[indigo, opacity=0.55, path fading=rglow] (8,4.5) circle (9);
  \end{scope}
  % aurora curtains
  \begin{scope}[blend mode=screen]
  \foreach \x/\c/\o in {2.5/ngreen/0.22, 4/nteal/0.18, 5.8/ncyan/0.20,
        9.5/nviolet/0.18, 11.5/nmag/0.16, 13.5/royal/0.18}{
     \fill[\c, opacity=\o, path fading=fadeup] (\x-0.7,4.5) rectangle (\x+0.7,9);}
  \end{scope}
  \dust{240}\bokeh{55}
  % the three repos as a constellation
  \coordinate (r1) at (4.6,6.3);\coordinate (r2) at (8.2,7.0);\coordinate (r3) at (11.4,5.6);
  \neon{(r1)}{(r2)}{ncyan}\neon{(r2)}{(r3)}{nmag}
  \spark{(r1)}{0.09}{ncyan}\spark{(r2)}{0.09}{npink}\spark{(r3)}{0.09}{nlime}
  \begin{scope}[blend mode=screen]
     \fill[ncyan, opacity=0.40, path fading=rglow] (8,3.2) ellipse (5.5 and 1.5);
  \end{scope}
  \node[text=white] at (8,3.3) {\fontsize{48}{52}\selectfont\textbf{Thank you}};
  \node[text=paper, opacity=0.7] at (8,2.1)
    {\footnotesize github.com/jaspek/rotor-rust \;\textcolor{nmag}{$\bullet$}\; rotor-c-hashcons
     \;\textcolor{nmag}{$\bullet$}\; jaspek.github.io/rotor-rust};
}

\end{document}
